% ---------------------------------------------------------------------------
% Chapter -- Feasibility Limits on Real Hardware
%
% DRAFT 1 (19 Oct 2026). Proposed by Claude, to be read and rewritten in
% your own words.
% EVERY number and every figure of this chapter comes from MATLAB:
%   "1 trailer/thesis_figures.m", "2 trailers/thesis_figures.m" and
%   thesis_figures_compare.m (they write figures/feas_*.pdf and feas_*.mat).
% The text quotes the printout of those functions; if you change params.m,
% run them again and compare the numbers below.
% ---------------------------------------------------------------------------

% a missing figure file prints a red box instead of stopping the compilation
\providecommand{\feasfig}[2]{%
  \IfFileExists{figures/#1.pdf}{\includegraphics[#2]{figures/#1}}%
  {\fbox{\textcolor{red}{figure \detokenize{#1} missing -- run thesis\_figures}}}}

\chapter{Feasibility Limits on Real Hardware}
\label{ch:feasibility}
\TODO[status of this chapter: about 80 percent done]

The controller of Chapter~\ref{ch:control} has a free gain $\gamma$ and a free
period $\varepsilon$, and the convergence it promises holds only for
sufficiently small $\varepsilon$. This chapter asks what these numbers can be
on the TurtleBot3, whose wheel speeds are limited. The answer turns out to
decide whether the method is usable at all: with the settings of the original
paper the commanded velocities are hundreds of times above the limits, and
once the controller is tuned to respect them, the car with two trailers needs
hours to reach the goal where the car with one trailer needs minutes.
Everything is computed with the MATLAB code described in
Chapter~\ref{ch:simulation}; the printouts behind the numbers are in
Appendix~\ref{app:feasibility}.

% ===========================================================================
\section{Problem formulation}
\label{sec:feas-problem}

\paragraph{Hardware limits.}
The TurtleBot3 Burger accepts a linear velocity of at most
\SI{0.22}{m/s} and an angular velocity of at most \SI{2.84}{rad/s}
\cite{robotis2026tb3}. In the notation of the thesis, with $u_1$ the driving
and $u_2$ the angular velocity,
\begin{equation}
  |u_1(t)|\le u_1^{\max}=\SI{0.22}{m/s},\qquad
  |u_2(t)|\le u_2^{\max}=\SI{2.84}{rad/s}.
  \label{eq:limits}
\end{equation}
The manufacturer quotes an operating time of about two and a half hours per
battery charge \cite{robotis2026tb3}; this matters in Section~\ref{sec:feas-two}.

\paragraph{An upper bound on the controls.}
The controller \eqref{eq:controller} is a sum of constant and oscillating
terms. By the triangle inequality its magnitude is at most the sum of the
magnitudes of the terms. For $r=3$,
\begin{equation}
  B_1=|a_1|+T_{12}+T_{112},\qquad
  B_2=|a_2|+T_{12}+2\,T_{112},
  \label{eq:bound3}
\end{equation}
and for $r=4$ the third oscillating term adds $2T_{1112}$ to $B_1$ and
$6T_{1112}$ to $B_2$, with $T_{1112}=(k_{1112}\omega)^{3/4}|a_{1112}|^{1/4}$.
The factors $2$ and $6$ are the amplitudes of the terms in the controller of
\cite{zuyev2024}. The bound for $u_1$ is not conservative: all cosines in $u_1$
equal one at the start of every interval, so $u_1$ attains $B_1$ there (up to
the sign of the small constant term $a_1$). The bound for $u_2$ is conservative,
since the sines vanish at the start of the interval and the peaks of the
different terms do not coincide; as the numbers below show, $u_1$ is the
limiting one anyway.

\paragraph{Three conditions on the tuning.}
A choice of $(\gamma,\varepsilon,k,V)$ is usable on the robot if

\begin{enumerate}
  \item[(i)] \emph{amplitude:} $B_1\le u_1^{\max}$ and $B_2\le u_2^{\max}$ at every
    point of the path the state follows;
  \item[(ii)] \emph{contraction:} $\varepsilon\gamma h(z)\le 0.8$ along the path, so that
    the factor $1-\mu$ in \eqref{eq:contraction} stays positive with a margin
    (the value $0.8$ is the safety margin used in the code, not a result of the theory);
  \item[(iii)] \emph{validity:} $\varepsilon$ is small enough for the nilpotent
    approximation to describe the real kinematics over one interval.
\end{enumerate}

Conditions (i) and (ii) are checked by computation (below). Condition (iii) is
not computed here: it is an observation from simulations of the exact kinematics
and is discussed in Chapter~\ref{ch:simulation}; where it matters we say so.

Write $\gamma_{\mathrm{amp}}(z)$ for the largest $\gamma$ satisfying (i) at
the state $z$. Because the amplitudes are proportional to $\gamma h(z)$, this is
found by a one-dimensional search. Along the straight ray from $z_0$ to the
origin, which every radial Lyapunov function follows (Chapter~\ref{ch:control}),
the largest admissible fixed gain is
\begin{equation}
  \gamma_{\mathrm{eff}}=\min\Bigl\{\min_{z\ \mathrm{on\ ray}}\gamma_{\mathrm{amp}}(z),\
  \frac{0.8}{\varepsilon\,\max_{\mathrm{ray}}h}\Bigr\}.
  \label{eq:geff}
\end{equation}
As the measure of performance we use the time $t_q$ for $\|z\|$ to drop to a
fraction $q$ of its initial value, computed from the gradient flow of
Section~\ref{sec:convergence}. It was checked against the full oscillating
simulation, with agreement to \SI{0.2}{\percent} of the progress over \SI{80}{s}
(\TODO[this check is quoted from the comments in the MATLAB code - repeat it for the final parameters in the simulation chapter]).
All numbers refer to the start pose of the experiment: for one trailer
$x_0=(-0.2\,\mathrm{m},-0.2\,\mathrm{m},45^\circ,45^\circ)$, with $\|z_0\|=0.835$; for two trailers
$x_0=0.6\,(-0.2\,\mathrm{m},-0.2\,\mathrm{m},45^\circ,45^\circ,45^\circ)=(-0.12\,\mathrm{m},-0.12\,\mathrm{m},27^\circ,27^\circ,27^\circ)$,
with $\|x_0\|=0.834$ and $\|z_0\|=0.502$. The start for two trailers was reduced by the factor $0.6$
so that $\|x_0\|<1$, as the supervisor required
(\TODO[confirm with the supervisor whether the condition $\|x_0\|<1$ refers to the physical state $x_0$, as assumed here, or to $z_0$; the one-trailer start has $\|x_0\|=1.15$]).
All numbers also refer to the measured
dimensions $d_0=\SI{115}{mm}$, $d_1=\SI{84}{mm}$, $\ell_1=\SI{82}{mm}$,
$\ell_2=\SI{89}{mm}$ (\TODO[the dimensions are those of the 15 September measurement - keep in sync with the parameter section of the model chapter]).

\paragraph{How the gain depends on the period.}
Every oscillating term of the controller has the form
$T_m=c\,(k\omega)^{(m-1)/m}|a_m|^{1/m}$ for a bracket of length $m$, and
$a_m$ is proportional to $\gamma$. If the term of the deepest bracket ($m=r$)
alone used up a limit $u^{\max}$, the condition would read
$c\,\omega^{(r-1)/r}\gamma^{1/r}\approx u^{\max}$, i.e.
\begin{equation}
  \gamma_{\mathrm{amp}}\ \propto\ \varepsilon^{\,r-1}\,(u^{\max})^{r}.
  \label{eq:scaling}
\end{equation}
Two consequences organize the whole chapter. A longer period $\varepsilon$ allows a
much larger gain, the more so the deeper the bracket (exponent $r-1$), and the limit
on the linear velocity enters to the power $r$. Both are confirmed by the numbers
below.

% ===========================================================================
\section{Two trailers: feasible, but only slowly}
\label{sec:feas-two}

\paragraph{The settings of the paper.}
The simulations of \cite{zuyev2024} use $\varepsilon=0.1$, $\gamma=1$,
multipliers $(8,6,5)$ and the start $z_0=(1,-2,2,-1,1)$, with
$d_0=d_1=0.1$ and $\ell_1=\ell_2=1$ (units are not stated). For these values
our code reproduces the printed coefficients
$\alpha=0.10495$, $\kappa=0.0024$, $\theta=0.0112$ (and $1,1.5,4$ for the second
parameter set of the paper), which is an independent check of the formulas of
Chapter~\ref{ch:nilpotent}. The bound \eqref{eq:bound3} for this setting is
$B_1\approx 132$ and $B_2\approx 349$, i.e.\ about $600$ times the limit of
$u_1$ if the units are metres and seconds. To respect the limits at these
$\varepsilon$ and $k$, the gain would have to be $\gamma\approx3\times10^{-11}$.
The paper does not address hardware limits (it gives only an order estimate of the
controls, with the constant omitted); the point is that its settings are of a different
order of magnitude from what the robot can execute.

\paragraph{What fits the robot.}
For the robot we choose (Chapter~\ref{ch:simulation}) $\varepsilon=\SI{3}{s}$,
multipliers $(14,5,1)$ and the square-root Lyapunov function with
$\delta=0.005$, whose factor $h$ is discussed in Section~\ref{sec:feas-V}. The
largest gain along the ray satisfying \eqref{eq:geff} is
$\gamma=3.50\times10^{-5}$. At $z_0$ this uses \SI{100}{\percent} of the $u_1$
limit and \SI{18}{\percent} of the $u_2$ limit, so $u_1$ is the binding one. The
breakdown of the bound, as a share of the limit, is:

\begin{table}[ht]
  \centering
  \caption{Share of the velocity limits used by each term of the bound at $z_0$,
  for the settings of the experiment.}
  \label{tab:terms}
  \begin{tabular}{lccccc}
    \toprule
     & $|a|$ & $T_{12}$ & $T_{112}$ & $T_{1112}$ & sum \\
    \midrule
    one trailer, $u_1$   & 0.002 & 0.133 & 0.707 & --    & 0.843 \\
    one trailer, $u_2$   & 0.001 & 0.010 & 0.110 & --    & 0.121 \\
    two trailers, $u_1$  & 0.000 & 0.073 & 0.504 & 0.423 & 1.000 \\
    two trailers, $u_2$  & 0.000 & 0.006 & 0.078 & 0.098 & 0.182 \\
    \bottomrule
  \end{tabular}
\end{table}

For the car with two trailers the two deepest brackets use \SI{93}{\percent} of the
linear-velocity budget, and the first-level oscillation and the constant part
together less than \SI{8}{\percent}. This is the quantitative form of the remark
of Chapter~\ref{ch:control} that the deep brackets are expensive: the gain is set
almost entirely by $T_{112}$ and $T_{1112}$. (The one-trailer row sums to $0.843$ at $z_0$ because with the logarithmic function
of the experiment, Section~\ref{sec:feas-V}, the amplitude $\propto h\,\|z\|$ is largest not at $z_0$
but at $\|z\|\approx0.3$ on the ray, and the gain is set there. In the simulation the peak of $|u_1|$
is $0.2141$ m/s ($97\%$ of the limit) on the exact kinematics and $0.2229$ m/s ($101.3\%$) on the approximation.
\TODO[the approximation exceeds the u1 limit by 1.3 percent on the actual path although the gain is the largest admissible along the ray; the path is not exactly the ray - decide whether to add a margin])

\paragraph{The price: time.}
With this gain the estimated times for $\|z\|$ to fall to \SI{50}{\percent},
\SI{10}{\percent} and \SI{0.1}{\percent} of $\|z_0\|$ are \SI{2.0}{h},
\SI{3.6}{h} and \SI{4.1}{h}. For the car with one trailer, with its own
settings ($\varepsilon=\SI{5}{s}$, $(k_{12},k_{112})=(3,2)$, the logarithmic
function with $c=10$ of the supervisor, $\gamma=2.1\times10^{-2}$), the estimates are \SI{4.5}{min} to \SI{10}{\percent}
and \SI{8.2}{min} to \SI{0.1}{\percent}; the simulation with the exact kinematics
reaches \SI{0.1}{\percent} after \SI{445}{s} (the approximate model also after
\SI{445}{s}). One additional trailer therefore costs a factor of about $30$ in time,
from minutes to hours. Two facts matter beyond the number itself:

\begin{itemize}
  \item \SI{4.1}{h} exceeds the manufacturer's operating time of the robot
    (\SI{2.5}{h}) \cite{robotis2026tb3} by a factor of $1.6$; a full convergence of the
    two-trailer system would need more than one battery charge, or a different
    objective (for instance stopping at \SI{10}{\percent}, which still takes about
    \SI{3.6}{h}).
  \item The \SI{4.1}{h} is an estimate from the gradient flow. The full oscillating
    simulation of the exact kinematics was compared with the approximation
    over the first minute (the distance between the two stays below
    \SI{2.7}{\percent} of $\|z_0\|$) and, interval by interval, over \SI{30}{min}
    (Chapter~\ref{ch:simulation}); a simulation over the whole time on the exact kinematics
    does not confirm the estimate
    (\TODO[a longer run on the exact kinematics at eps = 3 s shows that the norm of z stops decreasing after about 3450 s, at about 0.82 of its initial value, and then grows again; so the estimate of 4.1 h holds for the approximation only; decide how to present this here and in the simulation chapter]).
\end{itemize}

% ===========================================================================
\section{Which parameters help and which do not}
\label{sec:feas-geometry}

The geometry of the vehicle enters the controller in two ways: through the
coefficients $\alpha,\kappa,\theta$ of the nilpotent approximation, which multiply
the deep brackets in the amplitudes \eqref{eq:amps}, and through the coordinate
change $x\mapsto z$, which fixes how large the same physical offset looks in
$z$. We varied one dimension at a time while keeping the start pose
$x_0$ fixed (so $z_0$ changes with the geometry), the quadratic Lyapunov
function and the period of the experiment, and computed the time to reach
\SI{10}{\percent} of $\|z_0\|$.

\begin{figure}[ht]
  \centering
  \feasfig{feas_1tr_geom}{width=0.48\textwidth}\hfill
  \feasfig{feas_2tr_geom}{width=0.48\textwidth}
  \caption{Time to \SI{10}{\percent} of $\|z_0\|$ against the hitch offset $d_0$ (left in each
  pair) and a common scale factor on the trailer lengths (right in each pair); one trailer
  at $\varepsilon=\SI{5}{s}$ (minutes), two trailers at $\varepsilon=\SI{3}{s}$ (hours). Dotted lines
  mark the measured dimensions.}
  \label{fig:geometry-sweep}
\end{figure}

\begin{itemize}
  \item \emph{Hitch offset $d_0$} (from \SI{20}{mm} to \SI{300}{mm}). The effect is modest
    and has a different sign for the two systems: with one trailer the time falls
    from \SI{9.6}{min} to \SI{6.2}{min} as $d_0$ grows, with two trailers it grows from
    \SI{8.3}{h} to \SI{13.8}{h}.
  \item \emph{Trailer length $\ell$} (scale from $0.5$ to $4$). Longer trailers are slower:
    with one trailer from \SI{7.7}{min} ($0.5\times$) to \SI{23}{min} ($4\times$); with two
    trailers from \SI{7.3}{h} ($0.5\times$) through \SI{9.2}{h} ($1\times$) and
    \SI{19.5}{h} ($2\times$) to \SI{59}{h} ($4\times$).
\end{itemize}

For changes of the geometry one could realistically make on this platform the spread is a
factor of two to three, small next to the effect of the period, the multipliers
and the Lyapunov function (Sections~\ref{sec:feas-levers} and~\ref{sec:feas-V}).
The geometry is also not freely chosen: the measured dimensions are those of the
available hardware. We therefore regard $\ell$ as the only geometric parameter worth
adjusting (shorter hitches, if the hardware allows), and $d_0$ as essentially fixed.
(Why the two systems react to $d_0$ in opposite ways is not explained by this
computation; we do not claim an explanation.)

% ===========================================================================
\section{Comparison by degree of nonholonomy}
\label{sec:feas-r}

The three systems of the same family differ in the length $r$ of the longest bracket needed
to reach all directions: $r=2$ is the car alone (only $[f_1,f_2]$), $r=3$ the car with
one trailer, $r=4$ the car with two. \Cref{fig:cmp-eps} and \cref{tab:cmp} compare
them at equal period, with the quadratic $V$, the start pose of the experiment (for $r=2$ the first
three privileged coordinates of the one-trailer start; for $r=4$ the reduced start of
Section~\ref{sec:feas-two}, so the comparison is not at equal distance) and the same limits. The
case $r=2$ has no code of its own in the study; it is computed with the
same recipe from the terms $a_1,a_2,T_{12}$ only (see the header of
\texttt{thesis\_figures\_compare.m}), with $k_{12}=1$.

\begin{figure}[ht]
  \centering
  \feasfig{feas_cmp_eps}{width=0.95\textwidth}
  \caption{Left: the largest gain compatible with the velocity limits, $\gamma_{\mathrm{amp}}$,
  against the period $\varepsilon$, for $r=2,3,4$ (dotted: the contraction cap
  $0.8/\varepsilon$ for the quadratic $V$). Right: time to \SI{10}{\percent} of
  $\|z_0\|$ with $\gamma_{\mathrm{eff}}$. The shaded region marks where the
  two-trailer approximation was found to break down (\cref{sec:feas-problem}, (iii)).}
  \label{fig:cmp-eps}
\end{figure}

\begin{table}[ht]
  \centering
  \caption{Time to \SI{10}{\percent} of $\|z_0\|$ (quadratic $V$, limits of the TurtleBot3).}
  \label{tab:cmp}
  \begin{tabular}{rrrr}
    \toprule
    $\varepsilon$ [s] & $r=2$ [min] & $r=3$ [min] & $r=4$ [h]\\
    \midrule
    1 & 0.50 & 168 & 115\\
    2 & 0.28 & 46.6 & 22.9\\
    3 & 0.21 & 22.3 & 9.18\\
    4 & 0.19 & 13.4 & 4.87\\
    5 & 0.24 & 9.0 & 3.01\\
    8 & 0.38 & 4.0 & 1.12\\
    \bottomrule
  \end{tabular}
\end{table}

Three things can be read off. First, at the same period every further bracket level
costs one to two orders of magnitude in time: at $\varepsilon=\SI{3}{s}$ the car alone
needs about \SI{12}{s}, one trailer \SI{22}{min} and two trailers \SI{9.2}{h}. Second, the
slope of $\gamma_{\mathrm{amp}}(\varepsilon)$ on the log--log plot at small $\varepsilon$ is
$0.95$, $1.9$ and $2.5$ for $r=2,3,4$ (computed over $\varepsilon=0.25$--$\SI{1}{s}$), in line with the
exponent $r-1=1,2,3$ of \eqref{eq:scaling}; for $r=4$ the asymptotic value is approached
only slowly because the lower terms still matter. Third, the larger the period, the
smaller the gap between the systems, but the period is limited from above by condition
(iii): for two trailers the approximation tracks the exact kinematics up to about
$\varepsilon=\SI{4}{s}$ and breaks down at \SI{6}{s} \TODO[these two numbers are taken from comments in the MATLAB parameter file - repeat the comparison of the approximation with the exact model for the final geometry and report it in the simulation chapter]. In the
admissible range the two-trailer car is therefore bound to take about \SI{5}{h} or more with the quadratic $V$.

% ===========================================================================
\section{Quantifying the levers}
\label{sec:feas-levers}

Four quantities can be changed without changing the robot: the period $\varepsilon$,
the multipliers $k$, the Lyapunov function (next section) and the size of the
workspace, i.e.\ the starting distance. We take them in turn, with the quadratic $V$
and the settings of the experiment as reference.

\paragraph{Period.} See \cref{tab:cmp}: roughly, $t\propto\varepsilon^{-(r-1)}$ while
the amplitude limit is binding. The contraction cap takes over at about
$\varepsilon=\SI{20}{s}$ for $r=3$ and above \SI{32}{s} for $r=4$, far outside the validity range
of the approximation. For $r=3$ no breakdown has been seen for $\varepsilon\le\SI{7}{s}$
(\TODO[same check as above for one trailer]). The period is thus the strongest lever, and it is
capped by condition (iii), not by the hardware.

\paragraph{Multipliers.}
\Cref{tab:k} shows the largest gain at the period of the experiment, for several
admissible multipliers. The deepest bracket is the most expensive, so the
multiplier of the deepest term should be smallest ($T_m\propto k^{(m-1)/m}$): the
exponents are $1/2$, $2/3$ and $3/4$, and the shallow multiplier $k_{12}$ is almost free.
The paper's choice $(8,6,5)$ is $13$ times slower than $(14,5,1)$ for the same
Lyapunov function; $(18,4,1)$ is $22\%$ faster than $(14,5,1)$.
The multipliers cannot be shrunk without limit because of the non-resonance
condition (\cref{sec:multipliers}); $(14,5,1)$ and $(18,4,1)$ satisfy condition C1
of \cite{zuyev2024}. A larger $k_{12}$ also means a faster oscillation
($k_{12}\omega=\SI{29}{rad/s}$ for $14$, \SI{38}{rad/s} for $18$ at $\varepsilon=\SI{3}{s}$),
which the velocity controller of the robot has to follow
(\TODO[how well the velocity controller of the robot follows an oscillation of 4.7 to 6 Hz - the experiment will tell]).
The multipliers $(14,5,1)$ were chosen by the search of the MATLAB code; with the
current dimensions $(18,4,1)$ is better, probably because it lies outside the range
that was searched (\TODO[check the search range, then decide whether to retune to (18,4,1)]).

\begin{table}[ht]
  \centering
  \caption{Effect of the multipliers. Quadratic $V$: largest gain at $z_0$ and time to
  \SI{0.1}{\percent}. $V$ of the experiment (logarithmic with $c=10$ for one trailer, square root
  with $\delta=0.005$ for two): gain along the ray and time to \SI{0.1}{\percent}.}
  \label{tab:k}
  \begin{tabular}{lrrrr}
    \toprule
    $k$ & $\gamma_{\mathrm{amp}}$ & $t_{0.1\%}$, quadr. & $\gamma_{\mathrm{eff}}$, exper. & $t_{0.1\%}$, exper.\\
    \midrule
    \multicolumn{5}{l}{one trailer, $\varepsilon=\SI{5}{s}$ (times in min)}\\
    $(3,2)$ & $4.25\times10^{-3}$ & 27.1 & $2.11\times10^{-2}$ & 8.19\\
    $(5,2)$ & $3.73\times10^{-3}$ & 30.8 & $1.72\times10^{-2}$ & 10.1\\
    $(7,2)$ & $3.39\times10^{-3}$ & 34.0 & $1.47\times10^{-2}$ & 11.8\\
    $(2,3)$ & $2.39\times10^{-3}$ & 48.2 & $1.39\times10^{-2}$ & 12.5\\
    \midrule
    \multicolumn{5}{l}{two trailers, $\varepsilon=\SI{3}{s}$ (times in h)}\\
    $(14,5,1)$ & $6.97\times10^{-5}$ & 27.5 & $3.50\times10^{-5}$ & 4.06\\
    $(18,4,1)$ & $8.50\times10^{-5}$ & 22.6 & $4.27\times10^{-5}$ & 3.33\\
    $(8,6,5)$  & $5.32\times10^{-6}$ & 361 & $2.67\times10^{-6}$ & 53.3\\
    \bottomrule
  \end{tabular}
\end{table}

\paragraph{Size of the workspace.}
With the position part of the start scaled by a factor $s$ (the angles kept at
$45^\circ$), the time to \SI{10}{\percent} for $s=0.5,1,2,4$ is
\SI{3.3}{min}, \SI{9.0}{min}, \SI{27}{min}, \SI{91}{min} for one trailer and
\SI{4.6}{h}, \SI{9.2}{h}, \SI{22}{h}, \SI{64}{h} for two (here $s=1$ is the reduced start of
the experiment). Doubling the distance
multiplies the time by about $3$ (one trailer) or by $2.0$ to $2.9$ (two). Starting close to the
goal is therefore much cheaper, which is relevant for the experiment: the starting
distance of \SI{0.12}{m} in each coordinate (already reduced from \SI{0.2}{m}) still costs hours for two trailers.

\paragraph{Combined.}
The levers act together. Starting from $\varepsilon=\SI{3}{s}$, the paper's multipliers
and the quadratic $V$, the time to \SI{0.1}{\percent} for two trailers is
\SI{361}{h}; the choice $(14,5,1)$ brings it to \SI{27.5}{h}, the square-root $V$ to
\SI{4.1}{h}, and $(18,4,1)$ would give \SI{3.3}{h}. For one trailer, starting from
$\varepsilon=\SI{3}{s}$ with the quadratic $V$ ($\SI{67}{min}$), $\varepsilon=\SI{5}{s}$
gives \SI{27}{min} and the logarithmic $V$ of the experiment \SI{8.2}{min} (the square-root $V$
\SI{4.6}{min}). Each lever helps, but none changes the
qualitative picture: two trailers stay more than an order of magnitude slower than one.

% ===========================================================================
\section{State-dependent gain and the choice of Lyapunov function}
\label{sec:feas-V}

The amplitudes are proportional to $\gamma h(z)$ and, because $a_m\propto\|z\|$
at fixed $\gamma h$, the oscillation required shrinks as the state approaches the goal
($T_m\propto(\gamma h\|z\|)^{1/m}$). A constant gain, as with the quadratic $V$
($h=1$), is fixed by the worst point, which is the start; near the goal it leaves
most of the velocity budget unused. This suggests making the gain depend on the
state, and since $V$ enters only through $h$ (Chapter~\ref{ch:control}), this is the
same as changing the Lyapunov function. For a radial $V=\Phi(\|z\|)$, $h=\Phi'(r)/(\gamma r)$, and
the radius obeys $\dot r=-\gamma h(r)\,r$. We compare four choices, each with the
largest gain admissible along the ray by \eqref{eq:geff}:

\begin{itemize}
  \item quadratic: $h=1$;
  \item logarithmic: $V=\frac{\gamma}{2c}\ln(1+c\|z\|^2)$, $h=1/(1+c\|z\|^2)$ with $c=10$ and $c=100$ \cite{grushkovskaya2023};
  \item square root: $V=\gamma(\sqrt{\|z\|^2+\delta^2}-\delta)$, $h=1/\sqrt{\|z\|^2+\delta^2}$.
\end{itemize}

For the square root, $\dot r=-\gamma\,r/\sqrt{r^2+\delta^2}$: the speed is approximately
the constant $\gamma$ when $\|z\|\gg\delta$ and linear in $\|z\|$ for $\|z\|\ll\delta$. The
price is that $h\le1/\delta$, so the contraction condition (ii) becomes
$\gamma\le0.8\,\delta/\varepsilon$ and for small $\delta$ it is the cap, not the
amplitude, that eventually limits the gain. In the experiment the amplitude
limit is the binding one for both systems ($\varepsilon\gamma\,h_{\max}=0.11$ for one trailer with the
logarithmic function, $0.02$ for two trailers with the square-root function).

\begin{figure}[ht]
  \centering
  \feasfig{feas_1tr_V}{width=0.48\textwidth}\hfill
  \feasfig{feas_2tr_V}{width=0.48\textwidth}
  \caption{Estimated decay of $\|z\|/\|z_0\|$ for the four Lyapunov functions, each at its own
  largest admissible gain. Left: one trailer ($\varepsilon=\SI{5}{s}$, minutes);
  right: two trailers ($\varepsilon=\SI{3}{s}$, hours).}
  \label{fig:V}
\end{figure}

\begin{table}[ht]
  \centering
  \caption{Lyapunov functions at their largest admissible gain.
  ``Best radial'' is the lower limit for any radial $V$: the gain is chosen at every
  radius as large as the limits and the cap allow.}
  \label{tab:V}
  \begin{tabular}{lrrrr}
    \toprule
    $V$ & $\gamma_{\mathrm{eff}}$ & $t_{10\%}$ & $t_{0.1\%}$ \\
    \midrule
    \multicolumn{4}{l}{one trailer, $\varepsilon=\SI{5}{s}$ (min)}\\
    quadratic                    & $4.25\times10^{-3}$ & 9.0 & 27.1\\
    logarithmic, $c=10$          & $2.11\times10^{-2}$ & 4.5 & 8.2\\
    logarithmic, $c=100$         & $7.21\times10^{-2}$ & 8.5 & 9.7\\
    square root, $\delta=0.02$   & $3.20\times10^{-3}$ & 3.9 & 4.6\\
    best radial                  & --                  & 3.7 & 4.3\\
    \midrule
    \multicolumn{4}{l}{two trailers, $\varepsilon=\SI{3}{s}$ (h)}\\
    quadratic                    & $6.97\times10^{-5}$ & 9.2  & 27.5\\
    logarithmic, $c=10$          & $2.45\times10^{-4}$ & 4.0  & 9.3\\
    logarithmic, $c=100$         & $1.06\times10^{-3}$ & 3.9  & 5.1\\
    square root, $\delta=0.005$  & $3.50\times10^{-5}$ & 3.6  & 4.1\\
    best radial                  & --                  & 2.9  & 3.1\\
    \bottomrule
  \end{tabular}
\end{table}

Two observations. First, the state-dependent gain makes a large difference
to the \emph{late} part of the descent: to \SI{0.1}{\percent} the square-root
$V$ is $5.8$ times faster than the quadratic one with one trailer and $6.8$ times
faster with two. To \SI{10}{\percent} the difference is smaller ($2.3$ and $2.6$ times),
because the early part is limited by the same worst point. Second, the square-root
function is within \SI{7}{\percent} (one trailer) and \SI{30}{\percent} (two trailers) of the
best that any radial $V$ could reach under the same limits, so there is little to gain
from further searching in this family. The logarithmic function with $c=10$, which is used in the experiment
with one trailer, needs $8.2$ min to \SI{0.1}{\percent}, i.e.\ $1.8$ times as long as the square-root function and $1.9$ times the best radial
design (\TODO[state why the experiment uses the logarithmic function of the supervisor and not the faster square-root one]); the remaining gap in the two-trailer case is
the cap $\varepsilon\gamma h\le0.8$ with $\delta=0.005$.

A caveat on the theory: the convergence theorem of \cite{zuyev2024} is stated for a
quadratic $V$, and the logarithmic function comes from the related work
\cite{grushkovskaya2023}. The square-root function is smooth, positive definite, radial and
has a bounded gradient, but we have not checked it against the assumptions of the theorem
(\TODO[ask the supervisor which assumptions on V the convergence proof needs and whether the square-root function satisfies them]).
The comparison above is therefore a comparison of designs under the hardware limits, not
of proved guarantees.

% ===========================================================================
\section{Hardware requirements}
\label{sec:feas-hw}

What would a faster robot buy? We scaled each limit separately by a factor
$s$ (quadratic $V$, the period of the experiment). The result for the
time to \SI{10}{\percent}:

\begin{figure}[ht]
  \centering
  \feasfig{feas_1tr_hw}{width=0.48\textwidth}\hfill
  \feasfig{feas_2tr_hw}{width=0.48\textwidth}
  \caption{Time to \SI{10}{\percent} when one limit is scaled and the other kept at its
  TurtleBot3 value; left: one trailer, right: two trailers. The curve for $u_2$ is
  flat; the one for $u_1$ saturates when another constraint takes over (the cap for one
  trailer, the $u_2$ limit for two).}
  \label{fig:hw}
\end{figure}

\begin{itemize}
  \item \emph{Faster rotation does nothing.} Even at its present value $u_2$ is used to
    only \SI{12}{\percent} (one trailer) and \SI{18}{\percent} (two trailers) of its limit, so raising
    the limit on the angular velocity changes nothing.
  \item \emph{Faster driving helps a lot.} The time falls roughly as $(u_1^{\max})^{-r}$, as predicted by \eqref{eq:scaling}
    (measured exponents about $2.7$ for one trailer and $3.2$ for two, between
    $s=0.5$ and $s=2$). For two trailers, doubling the limit of $u_1$ to \SI{0.44}{m/s} gives
    \SI{1.0}{h}; $s=3$ (\SI{0.66}{m/s}) brings the time to \SI{10}{\percent} down to \SI{17}{min}. For one trailer,
    doubling gives \SI{1.4}{min}. The improvement stops at about $s=4$ for one trailer, where
    the cap $\varepsilon\gamma\le0.8$ binds, and at about $s=8$ for two trailers, where the
    limit of $u_2$ becomes the binding one (raising both limits together would help further).
  \item \emph{A larger vehicle} is not a remedy in itself. Scaling the trailers by a factor
    $2$ (\cref{sec:feas-geometry}) makes the two-trailer car slower by a factor of
    $2.1$, and a larger robot would also have larger velocity limits only if its drive
    allows it.
\end{itemize}

\paragraph{Summary.} For the TurtleBot3 the control of the car with one trailer is
feasible in a few minutes with the oscillating controller. For two trailers it is feasible in the
sense that conditions (i)--(iii) can all be met, but the price is an
estimated convergence time of the order of hours, set by the linear-velocity limit and the
deep brackets; the periods that would shorten it are excluded by the validity of the approximation. The
experiments of Chapter~\ref{ch:experiments} are therefore designed for one trailer, with
two trailers treated in simulation.
